% Lab 8 — SI Delay Analysis（原书 PDF p.91–102）
\bichapter{Lab 8：SI 延迟分析}{Lab 8: SI Delay Analysis}

\begin{objectiveBox}
完成本实验后，你应能够：
\begin{itemize}
  \item 修改常规 PrimeTime 运行脚本以启用 PrimeTime SI 串扰延迟（crosstalk delay）分析；
  \item 在 shell 与 GUI 中生成并解读 SI 报告，收集指导时序收敛的信息。
\end{itemize}
\end{objectiveBox}

\begin{durationBox}
约 \textbf{45} 分钟。
\end{durationBox}

\section{概述}
\label{lab8:overview}

\noindent\textit{Overview}

本实验涵盖：修改运行脚本以启用 SI 延迟分析、在 shell 中解读 SI 时序报告，以及在 GUI 中查看串扰直方图。

\subsection{相关文件与目录}
\label{lab8:files}

\noindent\textit{Relevant Files and Directories}

\begin{tabularx}{\textwidth}{@{}l X@{}}
\toprule
\textbf{路径 / 文件} & \textbf{说明} \\
\midrule
\cmd{ref/db/} & 工艺库 \\
\cmd{ref/design_data/} & 寄生参数与设计网表 \\
\cmd{ref/scripts/} & 支持脚本 \\
\cmd{ref/tcl_procs/} & 实用 Tcl 过程 \\
\cmd{lab8_si_delay/} & 本实验工作目录 \\
\quad\cmd{RUN.tcl} & 常规 PrimeTime 运行脚本 \\
\quad\cmd{SI-RUN.tcl} & PrimeTime SI 运行脚本 \\
\quad\cmd{.synopsys_pt.setup} & Setup 文件 \\
\quad\cmd{max_func_savesession/} & 常规 PrimeTime 已保存会话 \\
\quad\cmd{rpt/} & 运行脚本生成的报告 \\
\bottomrule
\end{tabularx}


\section{操作说明}
\label{lab8:instructions}

\noindent\textit{Instructions}

\subsection{任务 1：执行 PrimeTime SI 运行脚本}
\label{lab8:task1}

\noindent\textit{Task 1. Execute a Run Script for PrimeTime SI}

\begin{enumerate}
  \item 回答下列问题：
\begin{questionBox}
\textbf{问题 1.}\ 将“常规”PrimeTime 运行脚本修改为 PrimeTime SI 运行脚本，必须包含哪两个步骤？\\[0.35em]
\textbf{答：} 必须包含的两个步骤：
\begin{lstlisting}
# 在运行脚本开头设置该变量
set si_enable_analysis true
# 读寄生参数时保留耦合电容
read_parasitics -keep_coupling -format SBPF route_xtalk.sbpf
\end{lstlisting}
\end{questionBox}

  \item 进入目录 \cmd{lab8_si_delay}。可选：查看 SI 运行脚本 \cmd{SI-RUN.tcl} 及相关文件。

  \item 执行 PrimeTime SI：
\begin{lstlisting}
unix% pt_shell -f SI-RUN.tcl | tee -i sirun.log
\end{lstlisting}

  \item 运行脚本完成后，启动 PrimeTime 并恢复运行期间创建的已保存会话。
        会话目录名含时间戳，多次运行会保存在不同 Unix 目录中：
\begin{lstlisting}
unix% pt_shell
pt_shell> restore_session simax_func_savesession_<time_stamp>
\end{lstlisting}

\begin{noteBox}
PrimeTime 支持命令、选项、变量与文件名补全：键入少量字符后按 \textbf{Tab} 键。
\end{noteBox}

  \item 使用 Tcl 过程 \cmd{aa}（always ask）查找所有包含 \texttt{si\_} 的命令/变量：
\begin{lstlisting}
pt_shell> aa si_
\end{lstlisting}

\begin{noteBox}
上述 Tcl 过程可在 SolvNet 查阅，Doc Id \textbf{012959}。
\end{noteBox}
\end{enumerate}

\subsection{任务 2：在 Shell 中生成并解读 SI 报告}
\label{lab8:task2}

\noindent\textit{Task 2. Generate and Interpret SI Reports in the Shell}

\begin{enumerate}
  \item 报告本设计中的违例总数，并回答下列问题。

\begin{noteBox}
常规 PrimeTime 已通过时序；新出现的违例均与串扰相关。
常规 PrimeTime 运行的会话保存在 \cmd{max_func_savesession} 目录中（如有兴趣可对比）。
\end{noteBox}

\begin{questionBox}
\textbf{问题 2.}\ 本设计有多少个违例？\\[0.35em]
\textbf{答：} \
\begin{lstlisting}
pt_shell> report_analysis_coverage
\end{lstlisting}
共有 \textbf{14} 个违例，且均为 setup 检查。

\textbf{问题 3.}\ 本设计中 setup 的最坏 slack 有多大？\\[0.35em]
\textbf{答：} \
\begin{lstlisting}
pt_shell> report_analysis_coverage -status violated
\end{lstlisting}
最坏 setup slack 约为 \textbf{-0.1422}（\cmd{s4_op1_reg_30_/D}）。
\end{questionBox}

  \item 生成到达最大负 slack 端点的时序报告（建议复制粘贴端点名以免拼写错误）：
\begin{lstlisting}
pt_shell> report_timing -crosstalk -nets -transition \
-nosplit -to <endpoint>
\end{lstlisting}

\begin{noteBox}
可选：使用别名 \cmd{view} 在带滚动条的独立窗口中打开时序报告：
\begin{lstlisting}
pt_shell> view report_timing -crosstalk -nets -transition \
-nosplit -to <endpoint>
\end{lstlisting}
\cmd{view} 对应的 Tcl 过程见 SolvNet Doc Id \textbf{014947}。
若实验环境未安装 Tk 包中的 \cmd{wish} 可执行文件，\cmd{view} 别名将不可用。
\end{noteBox}

\begin{questionBox}
\textbf{问题 4.}\ 该路径上 delta delay 最大的网是哪条？最坏 delta delay 有多大？\\[0.35em]
\textbf{答：} 该路径上最大 delta delay 为 \textbf{0.1024\,ns}。

\textbf{问题 5.}\ 该 delta delay 如何表示（时间单位还是对 VDD 的比值）？\\[0.35em]
\textbf{答：} 以\textbf{时间单位}表示；主库时间单位为 ns。

\textbf{问题 6.}\ delta delay 为正还是为负？是否符合你的预期？\\[0.35em]
\textbf{答：} delta delay 为\textbf{正}。合理——这是 setup 报告，数据路径上应考虑最坏情况下的串扰\textbf{减速}效应。

\textbf{问题 7.}\ 与该大 delta delay 对应的 stage delay 是多少？\\[0.35em]
\textbf{答：} stage delay 为单元 + 网延迟：$0.5522 + 0.1025 = 0.6547$\,ns；\\
delta delay 与 stage delay 之比约为 \textbf{15\%}。

\textbf{问题 8.}\ 报告中 ``\texttt{\&}'' 符号的含义是什么？\\[0.35em]
\textbf{答：} 表示延迟使用详细网寄生数据计算。符号含义（摘自 \cmd{report_timing} man 页）：
\begin{center}
\small
\begin{tabular}{@{}ll@{}}
\toprule
\textbf{符号} & \textbf{反标类型} \\
\midrule
\texttt{H} & Hybrid annotation \\
\texttt{*} & SDF back-annotation \\
\texttt{\&} & RC network back-annotation \\
\texttt{\$} & RC pi back-annotation \\
\texttt{+} & Lumped RC \\
\texttt{<none>} & Wire-load model 或 none \\
\bottomrule
\end{tabular}
\end{center}
\end{questionBox}

  \item 对同一端点生成 hold 时序报告（也可用上箭头重复上一条命令并追加 \opt{-delay min}）：
\begin{lstlisting}
pt_shell> !! -delay min
\end{lstlisting}

\begin{questionBox}
\textbf{问题 9.}\ delta delay 为正还是为负？是否符合你的预期？\\[0.35em]
\textbf{答：} delta delay 现为\textbf{负}，合理——这是 hold 报告，数据路径上应考虑最坏情况下的串扰\textbf{加速}效应。
\end{questionBox}

\begin{noteBox}
注意：该 hold 路径上网延迟可能为负，但应关注的 stage delay（有意义延迟数值）为正。
\end{noteBox}

  \item 参照 Job Aid，查找并执行合适命令，找出 delta delay 最大的 victim 网：
\begin{questionBox}
\textbf{问题 10.}\ 本设计中最大的 delta delay 是多少？\\[0.35em]
\textbf{答：} 违例时序路径上 setup 与 hold 的最大 delta delay 均为 \textbf{0.1738\,ns}：
\begin{lstlisting}
pt_shell> report_si_bottleneck -cost_type delta_delay
\end{lstlisting}

\textbf{问题 11.}\ 报告的网是 victim 还是 aggressor？\\[0.35em]
\textbf{答：} 这些是 \textbf{victim} 网。

\textbf{问题 12.}\ 报告的网是否位于违例时序路径上？\\[0.35em]
\textbf{答：} 是。默认 \opt{-slack\_lesser\_than} 为 0，因此报告的 victim 网均位于 setup 或 hold 违例路径上。
\end{questionBox}

\begin{questionBox}
\textbf{问题 13.}\ 该命令的哪个开关可让 delta delay 较大的时钟网也出现在报告中（默认不包含）？\\
（提示：使用 man 页或 help。）\\[0.35em]
\textbf{答：} 开关为 \opt{-include\_clock\_nets}：
\begin{lstlisting}
# 键入下列内容（含短横线）后按 Tab 补全
pt_shell> report_si_bottleneck -
cost_type            max                  min                  nosplit
pre_commands         slack_lesser_than    include_clock_nets
max_nets             minimum_active_aggressors   post_commands
significant_digits
\end{lstlisting}
\end{questionBox}

\begin{lstlisting}
pt_shell> vman report_si_bottleneck
\end{lstlisting}

\begin{questionBox}
\textbf{问题 14.}\ 如何减小这些网上的 delta delay，从而消除或缓解时序违例？\\
（提示：参阅 Job Aid 或 \cmd{report_si_bottleneck} 的 man 页。）\\[0.35em]
\textbf{答：} 使用 \cmd{set_coupling_separation} 增大 victim 与 aggressor 间距，从而减小耦合电容。
\end{questionBox}

  \item 参照 Job Aid，查找并执行合适命令，找出 delta delay 与 stage delay 比值最大的 victim 网：
\begin{questionBox}
\textbf{问题 15.}\ 该报告中的 cost 表示什么（时间单位还是比值）？\\[0.35em]
\textbf{答：} cost 是 delta delay 与 stage delay 的\textbf{比值}：
\begin{lstlisting}
pt_shell> report_si_bottleneck -cost_type delta_delay_ratio
\end{lstlisting}

\textbf{问题 16.}\ 报告的网是 victim 还是 aggressor？\\[0.35em]
\textbf{答：} 报告的网是 \textbf{victim}。

\textbf{问题 17.}\ 如何减小这些网上的 delta delay？（提示同上。）\\[0.35em]
\textbf{答：} 大比值通常由 victim 转换较慢（弱 victim）引起；可通过\textbf{加大 victim 驱动}或在 victim 上\textbf{插入 buffer} 修复。
\end{questionBox}

  \item 参照 Job Aid，查找并执行合适命令，找出本设计中最强的 aggressor 网：
\begin{questionBox}
\textbf{问题 18.}\ 该报告中的 cost 表示什么（时间单位还是对 VDD 的比值）？\\[0.35em]
\textbf{答：} cost 是\textbf{相对 VDD 的比值}——该 aggressor 网在所有相关违例 victim 网上引起的切换 bump 之和：
\begin{lstlisting}
pt_shell> report_si_bottleneck -cost_type delay_bump_per_aggressor
\end{lstlisting}

\textbf{问题 19.}\ 报告的网是 victim 还是 aggressor？\\[0.35em]
\textbf{答：} 这些是 \textbf{aggressor} 网。

\textbf{问题 20.}\ aggressor 网本身是否一定位于违例时序路径上？\\[0.35em]
\textbf{答：} \textbf{不一定}。只有 victim 网必须位于违例时序路径上（由 \opt{-slack\_lesser\_than} 定义）。

\textbf{问题 21.}\ 如何减小这些 aggressor 网引起的 delta delay？（提示同上。）\\[0.35em]
\textbf{答：} \textbf{缩小}强 aggressor 驱动，或通过网隔离减小耦合电容。
\end{questionBox}

  \item 参照 Job Aid，查找并执行合适命令，报告最坏时钟偏斜及相关的串扰贡献：
\begin{questionBox}
\textbf{问题 22.}\ 该报告是否表明时钟网屏蔽良好？\\[0.35em]
\textbf{答：} \textbf{是}。最坏时钟偏斜未被 delta delay 贡献恶化：
\begin{lstlisting}
pt_shell> report_clock_timing -type skew -crosstalk -verbose
\end{lstlisting}
\end{questionBox}
\end{enumerate}

\subsection{任务 3：在 GUI 中生成 SI 报告}
\label{lab8:task3}

\noindent\textit{Task 3. Generate SI Reports in the GUI}

对于很长的时序报告，在报告中查找最大 delta delay 可能较困难；GUI 可简化该任务。

\begin{enumerate}
  \item 启动 GUI：
\begin{lstlisting}
pt_shell> start_gui
\end{lstlisting}

  \item 从 \textbf{Crosstalk} 下拉菜单选择 \textbf{Delta Delay Histogram}，在对话框中点击 \textbf{OK}。

  \item 在直方图中选择最右侧 bin，其中包含 setup 的最坏 delta delay。

  \item 从 \textbf{Crosstalk} 菜单选择 \textbf{Accumulated Bump Voltage Histogram}，点击 \textbf{OK}。
        该视图按 victim 网显示累积 bump 电压（从最好到最坏）。

  \item 探索 Crosstalk 菜单中的其他感兴趣选项。

  \item 通过 \textbf{File > Close GUI} 或 shell 命令 \cmd{stop_gui} 关闭 GUI。
\end{enumerate}

\noindent 恭喜——Lab~8 完成！
